\scanpage{37}
\[
 \mathscr S^\circ\longrightarrow\underline{\operatorname{Hom}}_{1/2\,\mathrm{add}}(\mathcal A,\mathrm{Ab})
\]
Nous ferons cette étude dans le
cas où $\mathscr S=\widehat\Delta=\mathrm{Ss}$, et dans
chacun des 2 cas suivants
\deleted{\unsure{en supposant qu'on a}}

$1^\circ$) $C$ est formé des groupes ab. $1/2$ simpliciaux
\deleted{\unsure{modulaires|multiplicatifs}}\inserted{\unsure{str.|ab.} $\partial$-parfaits}, $\mathcal A$ est la
% slanted
\margin{$C$ abélienne}
sous catégorie formée des
$\Phi(i,n)$ ($n\in\mathbb N^*$, \deleted{$\unsure{j\geq1}$} $i\geq1$) et $\Lambda^j\Psi_*$ ($j\geq0$).

On note que \inserted{pour $X\in\operatorname{Ob}\widehat\Delta$ fini,} la connaissance de
$\operatorname{Hom}_{\widehat\Delta}(X_{\unsure{\bullet|*}},M_*)$ comme foncteur
covariant en $M_*$ est réduit à
celle de $\operatorname{Hom}_{\widehat\Delta}(x_{\bullet},\deleted{\unsure{F}}\varphi)$ pour
$\varphi$ de $\mathcal A$. On peut encore
réduire \deleted{\unsure{aux}}\inserted{pour} $\varphi$ de $\mathcal A_0$ défini
ci-dessous --, et $=$ pour $\varphi$ parcourant
les $\Phi^i$.

$2^\circ$) $C_0$ \inserted{$\subset C$} \deleted{est} formé \unsure{:} des groupes ab
% slanted
\margin{$C_0$ est additive mais non abélienne.}
$1/2$ simpliciaux \deleted{\unsure{des}} spéciaux
(\deleted{\unsure{dans}}\inserted{i.e.} $\in C$ et les $\pi_i(M_{\bullet})$ sont
sans torsion), $\deleted{\unsure{\Phi}}\mathcal A_0\subset C_0$ sous cat.
de ceux qui sont de la forme
$\Phi^{\unsure{0|i}}=\Lambda^i\Phi_*\deleted{\unsure{\Psi}}$ ($i\geq1$) ou $\Psi^j=\Lambda^{\unsure{0|j}}\Psi_*$ ($j\geq0$).
% archive sheet number: 36
